%=====================================================================
% PART III OPENER
%=====================================================================
\thispagestyle{plain}
\structuralfigures{P3.}

\begin{center}
\begin{tcolorbox}[enhanced, width=0.92\textwidth, colback=greenhl!5,
  colframe=greenhl, boxrule=0pt, leftrule=5pt, arc=3pt, top=8pt, bottom=8pt]
\large\itshape\color{greenhl!60!black}
A virtual machine virtualizes the hardware, and pays for a whole operating
system to do it. Part~III asks a cheaper question: if the processes already
share one kernel, what would it take to stop them seeing each other? The answer
turns out to be two features Linux already had --- and an industry grew out of
combining them.
\end{tcolorbox}
\end{center}

\vspace{6pt}

\dsfigh{%
\begin{tikzpicture}[
  cbox/.style={rounded corners=3pt, draw=greenhl, line width=1pt, fill=greenhl!7,
               text=greenhl!55!black, font=\small\bfseries, align=center,
               minimum width=3.5cm, minimum height=1.0cm},
  hbox/.style={rounded corners=3pt, draw=greenhl, line width=1.8pt, fill=greenhl!16,
               text=greenhl!55!black, font=\small\bfseries, align=center,
               minimum width=3.5cm, minimum height=1.0cm},
  arr/.style={-{Stealth[length=2.4mm]}, greenhl!60, line width=1pt}
]
  \node[hbox] (c8)  at (0,0)    {Ch.\ 8\\Namespaces\\and cgroups};
  \node[cbox] (c9)  at (4.3,0)  {Ch.\ 9\\Docker and OCI};
  \node[cbox] (c10) at (8.6,0)  {Ch.\ 10\\Kubernetes};
  \node[cbox] (c11) at (12.9,0) {Ch.\ 11\\Lightweight\\and Secure};
  \draw[arr] (c8)  -- (c9);
  \draw[arr] (c9)  -- (c10);
  \draw[arr] (c10) -- (c11);
  \node[font=\scriptsize\itshape, text=greenhl!55!black, align=center, fill=white,
        inner sep=2pt] at (2.15,-1.1) {package it};
  \node[font=\scriptsize\itshape, text=greenhl!55!black, align=center, fill=white,
        inner sep=2pt] at (6.45,-1.1) {run a thousand};
  \node[font=\scriptsize\itshape, text=greenhl!55!black, align=center, fill=white,
        inner sep=2pt] at (10.75,-1.1) {when one kernel\\is not enough};
  \node[font=\scriptsize\itshape, text=accent, align=center] at (12.9,1.15)
       {the hypervisor returns};
\end{tikzpicture}}%
{Reading path through Part~III. The arc runs from kernel primitives to a cluster
scheduler --- and then back to Part~II, because Chapter~11's answer to container
isolation is to put a virtual machine underneath each one.}{part3map}

\newpage
